\documentclass[11pt]{article}
\usepackage{manual_docs_style}
\usepackage{adjustbox}
\usepackage{amsmath}
\usepackage{booktabs}
\usepackage{longtable}
\usepackage{array}
\usepackage{geometry}
\usepackage{graphicx}
\usepackage{makecell}
\usepackage{multirow}
\usepackage{placeins}
\usepackage{tabularx}

\newcolumntype{Y}{>{\centering\arraybackslash}X}
\newcolumntype{S}{>{\raggedright\arraybackslash}p{1.9cm}}
\newcolumntype{M}{>{\raggedright\arraybackslash}p{2.8cm}}

\begin{document}

\providecommand{\FloatBarrier}{}

\newcommand{\identityFieldRows}{%
\code{row_slug} & Question row slug \\
\code{agent} & Agent profile identifier \\
\code{desc_level} & Textual-description level \\
\code{task_output} & Task-output interface: code or direct \\
\code{noise_level} & Noise regime \\
\code{train_samples} & Number of labeled training examples \\
\code{simulator} & Simulator or environment name \\
}

\newcommand{\metricFamilyRows}{%
Top-1 accuracy & \code{average_top1_accuracy} & \code{std_per_sample}, \code{std_per_batch} & Mean top-1 accuracy; standard deviations over per-sample correctness and per-seed batch accuracy \\
Shortlist score & \code{average_shortlist_score} & none & Mean shortlist score \\
Cost & \code{avg_cost} & \code{std_cost} & Estimated cost \\
Cumulative tokens & \code{avg_cumulative_tokens} & \code{std_cumulative_tokens} & Cumulative token usage \\
Tool calls & \code{avg_tool_calls} & \code{std_tool_calls} & Tool-call count \\
Agent steps & \code{avg_num_steps} & \code{std_num_steps} & Agent step count \\
}

\newcommand{\consolidatedCsvSchemaTable}[1]{%
\begingroup
\small
\setlength{\tabcolsep}{4pt}
\renewcommand{\arraystretch}{1.08}
\begin{tabular}{@{}>{\raggedright\arraybackslash}p{0.30\textwidth}>{\raggedright\arraybackslash}p{0.62\textwidth}@{}}
\toprule
\textbf{Field} & \textbf{Description} \\
\midrule
#1
\bottomrule
\end{tabular}
\endgroup
}

\newcommand{\metricFamilyTable}[1]{%
\begingroup
\footnotesize
\setlength{\tabcolsep}{2pt}
\renewcommand{\arraystretch}{1.08}
\begin{tabular}{@{}>{\raggedright\arraybackslash}p{0.14\textwidth}>{\raggedright\arraybackslash}p{0.26\textwidth}>{\raggedright\arraybackslash}p{0.27\textwidth}>{\raggedright\arraybackslash}p{0.23\textwidth}@{}}
\toprule
\textbf{Metric} & \textbf{Mean/value field} & \textbf{Std field(s)} & \textbf{Description} \\
\midrule
#1
\bottomrule
\end{tabular}
\endgroup
}

\newcommand{\identityFieldsTable}{\consolidatedCsvSchemaTable{\identityFieldRows}}
\newcommand{\perRowSlugCsvMetricTable}{\metricFamilyTable{\metricFamilyRows}}
\newcommand{\perAgentCsvMetricTable}{\metricFamilyTable{\metricFamilyRows}}


\newcommand{\calibrationHeader}{%
\textbf{Experiment} & \textbf{D} & \textbf{Noise} & \textbf{Train} &
\textbf{Test} & \textbf{Row Slug} & \textbf{Mode} \\
}
\newcommand{\experimentHeader}{%
\textbf{Experiment} & \textbf{D} & \textbf{Noise} & \textbf{Train} &
\textbf{Test} & \textbf{Row Slug} & \textbf{Mode} \\
}
\newcommand{\calibrationRow}[7]{%
\code{#1} & \code{#2} & #3 & #4 & #5 & \code{#6} & \code{#7} \\
}
\newcommand{\experimentRow}[7]{%
\code{#1} & \code{#2} & #3 & #4 & #5 & \code{#6} & \code{#7} \\
}

\docTitle{Experiment Plan}
It describes how to use the previous stages and it creates the artifacts that are shown in the paper or used for human analysis.
All the results by default are saved in three folders:
\begin{itemize}
  \item \code{results/} This folder is committed to Git and contains results that are relatively small.
  \item \code{results_large_artifacts/} This folder is not committed to Git and contains anything that is large (agentic trajectories)
  \item \code{docs/overleaf_paper/generated} This folder contains any result that is shown in the paper or in the appendix
\end{itemize}
for the paper: tables and plots.

\phantomsection\label{docstart:experiment_plan}

\begingroup
\renewcommand{\arraystretch}{1.0}
\setlength{\tabcolsep}{2pt}
\tiny

\vspace{0.25em} 

\begin{longtable}{
  p{3.2cm}
  p{2.1cm}
  p{1.2cm}
  p{1.6cm}
  p{0.9cm}
  p{3.7cm}
  p{2.23cm}
}
\caption{Calibration experiment evaluated with \code{gpt-5.5}, \code{gemini-3.1-pro-preview}, \code{claude-opus-4.6}, and \code{minimax-m2.7}; models with reasoning-effort settings use high reasoning.}\label{tab:experiment-plan-calibration}\\

\toprule
\calibrationHeader
\midrule
\endfirsthead

\toprule
\calibrationHeader
\midrule
\endhead

\bottomrule
\endlastfoot

\calibrationRow{baseline_calibration}{high}{none}{0}{10}{frost_01234-anchor}{direct}

\end{longtable}

\vspace{0.8em}

\begin{longtable}{
  p{3.2cm}
  p{2.1cm}
  p{1.2cm}
  p{1.6cm}
  p{0.9cm}
  p{3.7cm}
  p{2.23cm}
}
\caption{Main paper experiments evaluated with \code{gpt-5.5}, \code{gemini-3.1-pro-preview}, \code{claude-opus-4-6}, and \code{minimax-m2.7}; models with reasoning-effort settings use high reasoning.}\label{tab:experiment-plan-main}\\

\toprule
\experimentHeader
\midrule
\endfirsthead

\toprule
\experimentHeader
\midrule
\endhead

\bottomrule
\endlastfoot

\experimentRow{baseline}{high}{low}{3}{10}{gentle_01234-cloud}{direct}
\experimentRow{code_baseline}{high}{low}{3}{10}{gentle_01234-flame}{code}
\experimentRow{baseline_noise_high}{high}{high}{3}{10}{gentle_01234-pine}{direct}
\experimentRow{code_baseline_noise_high}{high}{high}{3}{10}{gentle_01234-orbit}{code}

\end{longtable}

\vspace{0.8em}

\begin{longtable}{
  p{3.2cm}
  p{2.1cm}
  p{1.2cm}
  p{1.6cm}
  p{0.9cm}
  p{3.7cm}
  p{2.23cm}
}
\caption{Ablation experiments evaluated with \code{gpt-5.5} and \code{gemini-3.1-pro-preview}, both using high reasoning.}\label{tab:experiment-plan-ablation}\\

\toprule
\experimentHeader
\midrule
\endfirsthead

\toprule
\experimentHeader
\midrule
\endhead

\bottomrule
\endlastfoot

\experimentRow{direct_context_shot_3}{high}{high}{0}{10}{frost_01234-pine}{direct}
\experimentRow{code_context_shot_3}{high}{high}{0}{10}{frost_01234-orbit}{code}
\experimentRow{no_context_shot}{none}{high}{3}{10}{gentle_01234-harbor}{direct}
\experimentRow{code_no_context_shot}{none}{high}{3}{10}{gentle_01234-tide}{code}


\end{longtable}
\endgroup

\section{Consolidated Results}
These are two scripts that gather from each run folder all the results in \code{results_large_artifacts/*}.
The first script is:
\begin{DocCode}
python consolidate_trajectories.py [TAG ...] [--output-folder FOLDER]
\end{DocCode}
This script collects all the \code{atif_trajectory.json}, optionally filted by tag and them in \code{FOLDER} by default in \code{results_large_artifacts/trajectories}

\begin{DocCode}
  python consolidate_results.py [TAG ...] [--output-folder FOLDER]
\end{DocCode}
This script takes as input all \code{run_configurations/*.json}, optionally filters them by tag, and generate two files at \code{FOLDER} by default \code{results_large_artifacts/consolidated_numerical_results}
\code{results_large_artifacts/consolidated_numerical_results/<row_slug>_summary.json} and \code{results_large_artifacts/consolidated_numerical_results/<row_slug>.json}.
These two files have the following structure.
Both of them share the following identity fields:
\begingroup
\small
\setlength{\tabcolsep}{4pt}
\renewcommand{\arraystretch}{1.08}
\begin{tabular}{@{}>{\raggedright\arraybackslash}p{0.30\textwidth}>{\raggedright\arraybackslash}p{0.62\textwidth}@{}}
\toprule
\textbf{Field} & \textbf{Description} \\
\midrule
\code{timestamp} & \code{configuration_name.json::timestamp} \\
\code{configuration_name} & Name of the configuration file used \\
\code{paths} & Path to the scores \\
\termSimulator & \termSimulator name \\
\code{agent_id} & Agent profile identifier from \code{run_configurations/*.json::runs[]::agent_id} \\
\code{question_slug} & Question slug \\
\code{recipe} & Question recipe metadata copied from \code{question.recipe_info.*} field \\
\code{seed_numbers} & Int stating how many seeds have been collected with that specific question \\
\bottomrule
\end{tabular}
\endgroup


\code{<row_slug>_summary.json} contains additionally the following:
\begin{itemize}
  \item \code{score_batch}: a list of per-seed metric dictionaries from \code{scores.json::final_metric_test}. 
  Each dictionary contains \code{average_top1_accuracy}, \code{average_shortlist_score}, and \code{average_num_answers}.
  \item \code{score_other_batch}: a list of per-seed metric dictionaries from \code{scores.json::final_metric_other} (optional, only for code runs). Each dictionary contains \code{average_top1_accuracy}, \code{average_shortlist_score}, and \code{average_num_answers}.
  \item \code{trajectory_evaluation}: It is a list of dictionaries, where each dictionary has information directly from \code{trajectory_evaluation.json}
  \begin{itemize}
    \item \code{cost_usd}: \code{trajectory_evaluation.json::cost_usd}
    \item \code{tokens}: \code{trajectory_evaluation.json::tokens}
    \item \code{total_steps}:  \code{trajectory_evaluation.json::number_of_interactions.agent.steps}
    \item \code{number_of_tool_calls}: It is equal to \code{trajectory_evaluation.number_of_interactions.agent.tool_call_steps}
    \item \code{python_calls}: \code{trajectory_evaluation.json::subcommand_analysis.python.total_invocations} 
    \item \code{plots_generated}: number of entries in \code{trajectory_evaluation.json::artifact_analysis.plots} 
  \end{itemize}
\end{itemize}
While \code{<row_slug>.json} contains additionally the following:
\begin{itemize}
  \item \code{wrong_uuids}: It is a list of two-item lists, where each element contains \code{question_uuid} and \code{sample_uuid} for entries that the model answers incorrectly.
  \item \code{sample_results}: a list of per-run dictionaries copied from \code{scores.json::sample_results}.
\end{itemize}


In case there are multiple combinations of   \code{(agent_id, model, question_slug)} from different \code{run_configurations/*.json}, it returns the newest one based on the \code{configuration_name::<timestamp>} entry.
If any of these two files would be empty, then it's not created.

\section{Generate final artifacts}
\subsection{CSV Tables}
This script works with the output from all \code{<row_slug>.json}
\begin{DocCode}
  python generate_csv_tables.py [TAG ...]
\end{DocCode}
Where [TAG] if passed only includes results from the previous script with that corresponding tag.

The output of this second script consists of two CSV artifact types (for direct analysis and inspection)
\begin{itemize}
  \item \code{results/consolidated/all_runs.csv}: This is committed to Git
  \item \code{results/consolidated/tables/all_samples.csv}
\end{itemize}

These CSVs have the following columns

\paragraph{Run-level consolidated CSV.}
Each row in \code{results/consolidated/all_runs.csv} describes one consolidated run:

\allRunsCsvSchemaTable

\paragraph{Sample-level consolidated CSV.}
Each row in \code{results/consolidated/tables/all_samples.csv} describes one evaluated sample:

\allSamplesCsvSchemaTable

\FloatBarrier
\subsection{Generate TeX tables}
Use the following command:
\begin{DocCode}
  python generate_tex_tables.py
\end{DocCode}
This generates the final tables for the paper presented in the experimental section and appendix.
The input data comes from \code{all_run.csv}.
The tables are saved at \code{docs/overleaf_paper/generated/tables/main/*.tex} with name matching its label in the \code{.tex} and in the \code{used_by} entry.
So for example
\begin{DocCode}
\label{tab:acc_main_experiment_table}
\begingroup
\small
\newcommand{\scorepm}[2]{$#1\,{\scriptstyle\pm #2}$}

\setlength{\tabcolsep}{1.5pt}
\setlength{\aboverulesep}{0pt}
\setlength{\belowrulesep}{2pt}
\renewcommand{\arraystretch}{1.05}

\begin{tabular*}{\textwidth}{@{\extracolsep{\fill}}lcccccc@{}}
\toprule
& \multicolumn{3}{c}{\textbf{\texttt{Noise:} \texttt{low}}}
& \multicolumn{3}{c}{\textbf{\texttt{Noise:} \texttt{high}}} \\
\cmidrule(lr){2-4}\cmidrule(lr){5-7}
& \multicolumn{1}{c}{\texttt{Direct}}
& \multicolumn{2}{c}{\texttt{Code}}
& \multicolumn{1}{c}{\texttt{Direct}}
& \multicolumn{2}{c}{\texttt{Code}} \\
\cmidrule(lr){2-2}\cmidrule(lr){3-4}
\cmidrule(lr){5-5}\cmidrule(lr){6-7}
& \texttt{Episode}
& \texttt{Episode}
& \texttt{Held-out}
& \texttt{Episode}
& \texttt{Episode}
& \texttt{Held-out} \\
\midrule
\texttt{gpt-5.5}
& \scorepm{0.933}{0.09}
& \scorepm{0.853}{0.146}
& \scorepm{0.798}{0.106}
& \scorepm{0.833}{0.199}
& \scorepm{0.733}{0.238}
& \scorepm{0.665}{0.168} \\

\texttt{gemini-3.1-pro-preview}
& \scorepm{0.76}{0.164}
& \scorepm{0.733}{0.154}
& \scorepm{0.684}{0.139}
& \scorepm{0.56}{0.259}
& \scorepm{0.547}{0.196}
& \scorepm{0.57}{0.184} \\

\texttt{claude-opus-4-6}
& \scorepm{0.853}{0.125}
& \scorepm{0.687}{0.136}
& \scorepm{0.702}{0.133}
& \scorepm{0.64}{0.272}
& \scorepm{0.0}{0.0}
& \scorepm{0.0}{0.0} \\

\texttt{minimax-m2.7}
& \scorepm{0.293}{0.225}
& \scorepm{0.28}{0.193}
& \scorepm{0.353}{0.212}
& \scorepm{0.307}{0.227}
& \scorepm{0.364}{0.217}
& \scorepm{0.351}{0.218} \\
\bottomrule
\end{tabular*}
\endgroup

\end{DocCode}


In the following we show the tables as they appear in the paper.
Their numbers are correct; if a requested result is not present, the generated table renders \code{-}.
The model name is displayed in \code{shared/config/agents.json}.
\subsubsection*{Main experiment}
The data is from \code{baseline}, \code{code_baseline}, \code{baseline_noise_high}, and \code{code_baseline_noise_high} for \code{gpt-5.5}, \code{gemini-3.1-pro-preview}, \code{claude-opus-4.6} and \code{minimax-m2.7}
\par
\textbf{Accuracy table}\par
\begingroup
\small
\newcommand{\scorepm}[2]{$#1\,{\scriptstyle\pm #2}$}

\setlength{\tabcolsep}{3pt}
\setlength{\aboverulesep}{0pt}
\setlength{\belowrulesep}{2pt}
\renewcommand{\arraystretch}{1.05}

\begin{tabular*}{\textwidth}{@{\extracolsep{\fill}}lcccccc@{}}
\toprule
& \multicolumn{3}{c}{\textbf{\texttt{Noise:} \texttt{low}}}
& \multicolumn{3}{c}{\textbf{\texttt{Noise:} \texttt{high}}} \\
\cmidrule(lr){2-4}\cmidrule(lr){5-7}
& \multicolumn{1}{c}{\texttt{Direct}}
& \multicolumn{2}{c}{\texttt{Code}}
& \multicolumn{1}{c}{\texttt{Direct}}
& \multicolumn{2}{c}{\texttt{Code}} \\
\cmidrule(lr){2-2}\cmidrule(lr){3-4}
\cmidrule(lr){5-5}\cmidrule(lr){6-7}
& \texttt{Episode}
& \texttt{Episode}
& \texttt{Held-out}
& \texttt{Episode}
& \texttt{Episode}
& \texttt{Held-out} \\
\midrule
\texttt{gpt-5.5}
& \scorepm{0.96}{0.055}
& \scorepm{0.96}{0.055}
& \scorepm{0.888}{0.072}
& \scorepm{0.92}{0.084}
& \scorepm{0.86}{0.114}
& \scorepm{0.706}{0.092} \\

\texttt{gemini-3.1-pro-preview}
& \scorepm{0.82}{0.045}
& \scorepm{0.72}{0.303}
& \scorepm{0.686}{0.197}
& \scorepm{0.7}{0.158}
& \scorepm{0.42}{0.164}
& \scorepm{0.427}{0.035} \\

\texttt{claude-opus-4-6}
& \scorepm{0.88}{0.084}
& \scorepm{0.0}{0.0}
& \scorepm{0.0}{0.0}
& \scorepm{0.0}{0.0}
& \scorepm{0.0}{0.0}
& \scorepm{0.0}{0.0} \\

\texttt{minimax-m2.7}
& \scorepm{0.0}{0.0}
& \scorepm{0.0}{0.0}
& \scorepm{0.0}{0.0}
& \scorepm{0.0}{0.0}
& \scorepm{0.0}{0.0}
& \scorepm{0.0}{0.0} \\
\bottomrule
\end{tabular*}
\endgroup


\medskip
\textbf{Cost/token/tool-use main table template}\par
\begingroup
\small
\newcommand{\scorepm}[2]{$#1\,{\scriptstyle\pm #2}$}

\setlength{\tabcolsep}{5pt}
\setlength{\aboverulesep}{0pt}
\setlength{\belowrulesep}{2pt}
\renewcommand{\arraystretch}{1.05}

\begin{tabular*}{\textwidth}{@{\extracolsep{\fill}}lcccc@{}}
\toprule
& \multicolumn{2}{c}{\textbf{\texttt{Noise:} \texttt{low}}}
& \multicolumn{2}{c}{\textbf{\texttt{Noise:} \texttt{high}}} \\
\cmidrule(lr){2-3}\cmidrule(lr){4-5}
& \texttt{Direct}
& \texttt{Code}
& \texttt{Direct}
& \texttt{Code} \\
\midrule
\texttt{gpt-5.5} & \scorepm{0.559}{0.142} & \scorepm{0.664}{0.229} & \scorepm{0.7}{0.23} & \scorepm{0.714}{0.085} \\
\texttt{gemini-3.1-pro-preview} & \scorepm{1.088}{0.257} & \scorepm{1.839}{0.77} & \scorepm{1.282}{0.352} & \scorepm{1.542}{0.308} \\
\texttt{claude-opus-4-6} & \scorepm{2.53}{1.215} & \scorepm{0.0}{0.0} & \scorepm{0.0}{0.0} & \scorepm{0.0}{0.0} \\
\texttt{minimax-m2.7} & \scorepm{0.0}{0.0} & \scorepm{0.0}{0.0} & \scorepm{0.0}{0.0} & \scorepm{0.0}{0.0} \\
\bottomrule
\end{tabular*}
\endgroup


\subsubsection*{Ablation experiment}
The data is from \code{direct_context_shot_3}, \code{code_context_shot_3}, \code{no_context_shot}, and \code{code_no_context_shot} only for  \code{gpt-5.5} and \code{gemini-3.1-pro-preview}
\medskip
\par
\textbf{Accuracy-style ablation table template}\par
\begingroup
\small
\newcommand{\scorepm}[2]{$#1\,{\scriptstyle\pm #2}$}

\setlength{\tabcolsep}{5pt}
\setlength{\aboverulesep}{0pt}
\setlength{\belowrulesep}{2pt}
\renewcommand{\arraystretch}{1.05}

\resizebox{\textwidth}{!}{%
\begin{tabular}{lcccccc}
\toprule
& \multicolumn{3}{c}{Zero-example, high textual context}
& \multicolumn{3}{c}{Three-example, no textual context} \\
\cmidrule(lr){2-4}\cmidrule(lr){5-7}
& \multicolumn{1}{c}{\texttt{Direct}}
& \multicolumn{2}{c}{\texttt{Code}}
& \multicolumn{1}{c}{\texttt{Direct}}
& \multicolumn{2}{c}{\texttt{Code}} \\
\cmidrule(lr){2-2}\cmidrule(lr){3-4}
\cmidrule(lr){5-5}\cmidrule(lr){6-7}
& \texttt{Episode}
& \texttt{Episode}
& \texttt{Held-out}
& \texttt{Episode}
& \texttt{Episode}
& \texttt{Held-out} \\
\midrule
\texttt{gpt-5.5}
& \scorepm{.97}{0.100}
& \scorepm{.93}{0.100}
& \scorepm{.88}{0.100}
& \scorepm{.97}{0.100}
& \scorepm{.93}{0.100}
& \scorepm{.88}{0.100} \\

\texttt{gpt-5.3-codex-spark}
& \scorepm{.85}{0.100}
& \scorepm{.91}{0.100}
& \scorepm{.82}{0.100}
& \scorepm{.85}{0.100}
& \scorepm{.91}{0.100}
& \scorepm{.82}{0.100} \\

\texttt{gemini-3.1-pro-preview}
& \scorepm{.87}{0.100}
& \scorepm{.79}{0.100}
& \scorepm{.94}{0.100}
& \scorepm{.87}{0.100}
& \scorepm{.79}{0.100}
& \scorepm{.94}{0.100} \\

\texttt{claude-opus-4-6}
& \scorepm{.86}{0.100}
& \scorepm{.89}{0.100}
& \scorepm{.83}{0.100}
& \scorepm{.86}{0.100}
& \scorepm{.89}{0.100}
& \scorepm{.83}{0.100} \\

\bottomrule
\end{tabular}%
}
\endgroup


\medskip
\textbf{Cost/token/tool-use ablation table template}\par
\begingroup
\small
\newcommand{\scorepm}[2]{$#1\,{\scriptstyle\pm #2}$}

\setlength{\tabcolsep}{5pt}
\setlength{\aboverulesep}{0pt}
\setlength{\belowrulesep}{2pt}
\renewcommand{\arraystretch}{1.05}

\resizebox{\textwidth}{!}{%
\begin{tabular}{lcccc}
\toprule
& \multicolumn{2}{c}{Zero-example, high context}
& \multicolumn{2}{c}{Three-example, no context} \\
\cmidrule(lr){2-3}\cmidrule(lr){4-5}
& \texttt{Direct}
& \texttt{Code}
& \texttt{Direct}
& \texttt{Code} \\
\midrule
\texttt{gpt-5.5 (high)} & \scorepm{0.0}{0.0} & \scorepm{0.0}{0.0} & \scorepm{0.0}{0.0} & \scorepm{0.0}{0.0} \\
\bottomrule
\end{tabular}%
}
\endgroup


\clearpage
\textbf{Per-row-slug appendix template}\par
These tables are repeated for each experiment row in Tables~\ref{tab:experiment-plan-calibration}, \ref{tab:experiment-plan-main}, and \ref{tab:experiment-plan-ablation}.

\begingroup
\small
\newcommand{\scorepm}[2]{$#1\,{\scriptstyle\pm #2}$}

\setlength{\tabcolsep}{5pt}
\setlength{\aboverulesep}{0pt}
\setlength{\belowrulesep}{2pt}
\renewcommand{\arraystretch}{1.05}

\begin{tabular*}{\textwidth}{@{\extracolsep{\fill}}lccc@{}}
\toprule
& \textbf{\texttt{BallDrop}}
& \textbf{\texttt{DampedMass}}
& \textbf{\texttt{Inclined Plane}} \\
\midrule
\texttt{gpt-5.5}
& \scorepm{.97}{0.100}
& \scorepm{.88}{0.100}
& \scorepm{.93}{0.100} \\

\texttt{gemini-3.1-pro-preview}
& \scorepm{.85}{0.100}
& \scorepm{.82}{0.100}
& \scorepm{.91}{0.100} \\

\texttt{claude-opus-4-6}
& \scorepm{.87}{0.100}
& \scorepm{.94}{0.100}
& \scorepm{.79}{0.100} \\

\texttt{minimax-m2.7}
& \scorepm{.86}{0.100}
& \scorepm{.83}{0.100}
& \scorepm{.89}{0.100} \\

\bottomrule
\end{tabular*}
\endgroup


The tables are generated with the following stylistic rules:
\begin{itemize}
\item The standard deviation is the \code{standard deviation per batch}
\item Each number should display at most four digits, i.e., 0.559 is OK, 0.5593 is not OK; 123.4 is OK, 123.43 is not OK.

\end{itemize}

\begin{itemize}
  \item \code{*_main_experiment_table.tex} gets the data from the following experiments \code{baseline}, \code{code_baseline}, \code{baseline_noise_high}, and \code{code_baseline_noise_high}.
  \item \code{*_ablation_training_examples.tex} gets the data from the following experiments \code{direct_context_shot_3}, \code{code_context_shot_3}, \code{no_context_shot}, and \code{code_no_context_shot}.
\end{itemize}




\paragraph{Metric prefixes.}
Generated aggregate tables use these metric prefixes:

\begin{itemize}
\item \code{acc_} : accuracy 
\item \code{cost_}: cost 
\item \code{cum_}: cumulative tokens (in M) (so numbers must be divided by \code{1_000_000})
\item \code{num_steps}: num steps
\item \code{python_calls}: num of python calls
\item \code{plots_generated}: num of plots generated
\end{itemize}

\subsection{Plots}
This script generates the plots in the experiment section. 
It currently generates a figure with two subplots: (i) Accuracy vs Cumulative Tokens and (ii) Accuracy vs Python Calls
with only the data from the main experiments
\begin{DocCode}
  python generate_plots.py 
\end{DocCode}
The figure being generated is the following and it includes only the results in the main 

\begin{figure}[ht]
\centering
\includegraphics[width=\linewidth]{overleaf_paper/assets/figures/results/accuracy_vs_tokens_and_python_calls.pdf}
\caption{Top-1 accuracy as a function of cumulative token usage and Python calls for the results in the main experiment table.}
\label{fig:experiment-plan-accuracy-vs-compute}
\end{figure}

\end{document}
